\documentclass{article}
\usepackage{amsmath, amssymb, amsthm}
\usepackage{hyperref}
\usepackage{geometry}
\geometry{margin=1in}

\title{Erd\H{o}s Problem \#955}
\date{Last updated: 2025-08-31}
\author{Source: erdosproblems.com}

\begin{document}
\maketitle

\noindent\textbf{Status:} OPEN \quad \textbf{No prize}

\noindent\textbf{Tags:} number theory

\medskip
\noindent\textbf{Problem Statement:}

Let\[s(n)=\sigma(n)-n=\sum_{\substack{d\mid n\\ d<n}}d\]be the sum of proper divisors function.

If $A\subset \mathbb{N}$ has density $0$ then $s^{-1}(A)$ must also have density $0$.

\bigskip
\noindent\textbf{Remarks:}

A conjecture of Erd\H{o}s, Granville, Pomerance, and Spiro \cite{EGPS90}. It is possible for $s(A)$ to have positive density even if $A$ has zero density (for example taking $A$ to be the product of two distinct primes). Erd\H{o}s \cite{Er73b} proved that there are sets $A$ of positive density such that $s^{-1}(A)$ is empty.

Pollack \cite{Po14b} has shown that this is true if $A$ is the set of primes. Troupe \cite{Tr15} has shown that this is true if $A$ is the set of integers with unusually many prime factors. Troupe \cite{Tr20} has also shown this is true if $A$ is the set of integers which are the sum of two squares.

Pollack, Pomerance, and Thompson \cite{PPT18} prove that if $\epsilon(x)=o(1)$ and $A\subset \mathbb{N}$ has size at most $x^{1/2+\epsilon(x)}$ then\[\#\{ n\leq x: s(n)\in A\} =o(x)\]as $x\to \infty$. It follows that (using $s(n)\ll n\log\log n$) if $A$ grows like $\lvert A\cap [1,x]\rvert\leq x^{1/2+o(1)}$ then $s^{-1}(A)$ has density $0$.

Alanen calls $k$ such that $s(n)=k$ has no solutions untouchable. These are discussed in problem B10 of Guy's collection \cite{Gu04}.

\bigskip
\noindent\textbf{References:}

[EGPS90] Erd\H{o}s, P. and Granville, A. and Pomerance, C. and Spiro, C., On the normal behavior of the iterates of some arithmetic functions. Analytic number theory (Allerton Park, IL, 1989) (1990), 165-204.
 
 [Er73b] Erd\H{o}s, P., \"{U}ber die Zahlen der Form $\sigma (n)-n$ und $n-\phi(n)$. Elem. Math. (1973), 83-86.
 
 [Gu04] Guy, Richard K., Unsolved problems in number theory. (2004), xviii+437.
 
 [PPT18] Pollack, Paul and Pomerance, Carl and Thompson, Lola, Divisor-sum fibers. Mathematika (2018), 330--342.
 
 [Po14b] Pollack, Paul, Some arithmetic properties of the sum of proper divisors and
the sum of prime divisors. Illinois J. Math. (2014), 125--147.
 
 [Tr15] Troupe, Lee, On the number of prime factors of values of the
sum-of-proper-divisors function. J. Number Theory (2015), 120--135.
 
 [Tr20] Troupe, Lee, Divisor sums representable as the sum of two squares. Proc. Amer. Math. Soc. (2020), 4189--4202.

\bigskip
\noindent\small{Source: \url{https://www.erdosproblems.com/955}}
\end{document}
